\documentclass[11pt,a4paper]{article}
\usepackage[margin=2.4cm]{geometry}
\usepackage[T1]{fontenc}
\usepackage{lmodern}
\usepackage{amsmath,amssymb}
\usepackage{xcolor}
\usepackage{mdframed}
\usepackage{enumitem}
\usepackage{booktabs}
\usepackage{hyperref}
\hypersetup{colorlinks=true,linkcolor=blue!55!black,urlcolor=blue!55!black}
\usepackage{titlesec}

\definecolor{revblue}{RGB}{20,70,140}
\definecolor{quotegray}{RGB}{245,245,247}
\definecolor{newtext}{RGB}{0,90,50}

\titleformat{\section}{\large\bfseries\color{revblue}}{\thesection.}{0.6em}{}
\titleformat{\subsection}{\normalsize\bfseries}{\thesubsection}{0.6em}{}

% reviewer comment box
\newenvironment{comment}
  {\begin{mdframed}[backgroundcolor=quotegray,linewidth=0pt,skipabove=8pt,
    skipbelow=6pt,innertopmargin=6pt,innerbottommargin=6pt]\itshape}
  {\end{mdframed}}

% verbatim-ish quotation of new manuscript text
\newenvironment{newtextbox}
  {\begin{mdframed}[linecolor=newtext!45,linewidth=0.8pt,skipabove=6pt,
    skipbelow=8pt,innertopmargin=6pt,innerbottommargin=6pt]\small}
  {\end{mdframed}}

\newcommand{\where}[1]{\textcolor{revblue}{\textbf{[#1]}}}

\setlength{\parskip}{0.45em}
\setlength{\parindent}{0pt}

\begin{document}

\begin{center}
{\Large\bfseries Response to Reviewers}\\[0.6em]
{\large Manuscript ID: ca5f4c68-82d0-46f5-99ec-5304715a69a6}\\[0.4em]
{\large\itshape Discover Artificial Intelligence}\\[0.2em]
{\itshape Collection: Trustworthy and Responsible Federated Learning}
\end{center}

\vspace{0.5em}

\begin{tabular}{@{}p{3.1cm}p{13.2cm}@{}}
\textbf{Revised title} & Unifying Fairness Privacy Robustness and Explainability
in Federated Graph Neural Networks for High Stakes Risk Detection \\[0.35em]
\textbf{Previous title} & Trustworthy Federated Graph Neural Networks: Unifying
Fairness, Privacy, Robustness, and Explainability for High-Stakes Risk
Detection \\[0.35em]
\textbf{Decision} & Minor revision \\[0.35em]
\textbf{Enclosed} & Clean revised manuscript; tracked-changes manuscript
(all edits marked); this response letter \\
\end{tabular}

\vspace{1em}

Dear Dr.\ Fisichella and Ms.\ Baig, and dear Reviewers,

We thank the Editor and both Reviewers for a careful, constructive and unusually
useful set of comments. We are grateful that the reviewers found the work
technically sound and well motivated, and we have taken the criticisms seriously:
\textbf{every point raised has been addressed}, and in several cases we went
further than asked.

Three of the comments identified places where our text claimed more than our own
tables supported, and we want to acknowledge that directly rather than
defensively. Reviewer~1 was right that the ``only aggregator that stays strong
under all three attacks'' claim is contradicted by our own Table~9, right that
one baseline was counted but never reported, and right that five numbers
disagreed between text and tables. We have corrected all of them, and we also
re-audited \emph{every} numerical claim in the paper against the tables it cites,
which surfaced four further inconsistencies the reviewers did not raise. Those
are listed in Section~\ref{sec:extra} of this letter, in the interest of a clean
record.

In this letter, reviewer text appears in \emph{grey boxes}, our response follows
in plain text, and new or rewritten manuscript text appears in
\textcolor{newtext}{green-bordered boxes}. Page numbers refer to the clean
revised manuscript. Nothing in the revision required new experiments: as
Reviewer~1 anticipated, all changes concern wording, scope and consistency.
Two figures/tables were re-rendered from the \emph{existing} logged results
(no re-training) where the presentation itself was at fault; we flag both
explicitly below.

\vspace{0.5em}
\hrule
\vspace{0.5em}

\section{Summary of changes}

\begin{tabular}{@{}p{1.6cm}p{7.4cm}p{7.0cm}@{}}
\toprule
\textbf{Raised by} & \textbf{Issue} & \textbf{Where fixed} \\
\midrule
Editor & Title contains punctuation & New title (title page) \\
Editor & Declarations need separate subheadings & Declarations, p.~37 \\
Editor & Data availability statement & Declarations, p.~37 \\
Editor/R1 & Overstated claims & \S5.8 p.~26, \S5.2 p.~19, \S6 p.~29 \\
Editor/R1 & Text/table consistency & \S5.2, \S5.3, \S5.5, \S5.8 \\
Editor/R1 & Scope of the DP guarantee & Abstract, \S1, Table~1 note, \S4.2 \\
Editor/R2 & Protected attribute at inference & New paragraph, \S6 p.~32 \\
Editor/R2 & Trust-score limitations & \S6 p.~29 (moved before the score) \\
R1.1 & \texttt{robust\_bfwa} exclusivity claim & \S5.8 p.~26, \S5.3 p.~20 \\
R1.2 & Sixteen vs.\ fifteen baselines & Throughout; \S3.2 p.~15 \\
R1.3 & Uneven baseline coverage & New paragraph, \S3.2 p.~17 \\
R1.4 & Single-seed robustness & New paragraph \S5.8 p.~26; Limitation (4) \\
R1.5 & Energy column reads 0.00\,Wh & Table~11 and Table~10, \S5.10 \\
R1.6 & Five text/table mismatches & \S5.2, \S5.3, \S5.5 \\
R1.7 & DP scope in the abstract & Abstract; Table~1 note \\
R2.1 & Inference-time $s$ needs a full treatment & New paragraph \S6 p.~32; Limitation (12) \\
R2.2 & Trust-score caveat comes too late & \S6 p.~29; Limitation (10) \\
\midrule
\textit{Authors} & Four further self-identified inconsistencies & Section~\ref{sec:extra} below \\
\bottomrule
\end{tabular}

\vspace{0.8em}

\section{Response to the Editor}

\subsection*{E1. Revised title}

\begin{comment}
Revise Title: To aid our readers, and to maximize the accessibility of your
manuscript, the title should have a clear, precise scientific meaning. Where
possible, the title should be read as one concise sentence, without the use of
punctuation (full stops, colons, hyphens, question marks etc.). Please could you
re-write the title ensuring that it is informative and appropriate?
\end{comment}

Done. The previous title used a colon, three commas and a hyphen. The revised
title contains \textbf{no punctuation of any kind}, reads as a single sentence,
and still names all four trust properties and the application domain:

\begin{newtextbox}
\textbf{Unifying Fairness Privacy Robustness and Explainability in Federated
Graph Neural Networks for High Stakes Risk Detection}
\end{newtextbox}

We removed the hyphen from ``High-Stakes'' throughout the title and kept a short
running head (``Trustworthy Federated Graph Neural Networks'') for the page
footer, as the journal template requires.

\subsection*{E2. Declarations section with individual subheadings}

\begin{comment}
Authors are requested to include a separate section in the manuscript titled
``Declarations.'' In this section, please clearly state the following items as
individual subheadings: Ethical approval; Consent to participate; Consent to
publish. Each item should be addressed separately and appropriately.
\end{comment}

Done \where{Declarations, p.~37}. The Declarations were previously a single
bulleted list. They are now a section in which \emph{each} item carries its own
subheading: Funding; Conflict of interest; \textbf{Ethical approval};
\textbf{Consent to participate}; \textbf{Consent to publish}; Data availability;
Materials availability; Code availability; Author contributions. Each is
addressed substantively rather than with a bare ``not applicable'' --- for
example:

\begin{newtextbox}
\textbf{Ethical approval.} Not applicable. This study involved no human
participants and no animal subjects. It uses only publicly released,
de-identified benchmark datasets (German Credit, Credit Defaulter,
Bail/Recidivism, Pokec, Elliptic, and ogbn-products), each obtained from its
published distribution and used in accordance with its licence and intended
research purpose. No ethics-committee approval was therefore required.
\end{newtextbox}

\subsection*{E3. Data availability statement}

\begin{comment}
When submitting your revised manuscript, please select `yes' under the data
availability declaration $\dots$ Please make sure that data availability
statement in the manuscript and system exactly matches.
\end{comment}

Done \where{Declarations, p.~37}. The manuscript now contains an explicit
\textbf{Data availability} subheading naming every dataset and its public
source with URLs, and stating that no new data were generated. We will select
``yes'' in the submission system and paste this statement verbatim so that the
two match exactly.

We also took the opportunity to strengthen the \textbf{Code availability}
statement. The previous wording (``available from the corresponding author upon
request'') sat awkwardly beside the paper's reproducibility claims, so the
repository is now cited directly:
\url{https://github.com/vinhqdang/FedFairGNN}.

\subsection*{E4. Overstated claims, text/table consistency, DP scope}

\begin{comment}
For the revision, please address the few overstated or unsupported claims,
carefully check the consistency between the text and reported tables, and
clarify the scope of the differential privacy guarantee.
\end{comment}

These three requests coincide with Reviewer~1's comments 1, 6 and 7, and we
address them in full there (Sections~\ref{sec:r1-1}, \ref{sec:r1-6},
\ref{sec:r1-7} of this letter). In summary:

\begin{itemize}[leftmargin=1.2em,itemsep=2pt]
\item \textbf{Overstated claims.} The \texttt{robust\_bfwa} exclusivity claim is
withdrawn and replaced with an accurate two-aggregator statement (\S5.8); the
claim that all non-private baselines significantly beat us on German AUC is
corrected to name the exception (\S5.2); the ``lowest DPD/EOD of any method''
claim for PUFFLE is corrected to a tie (\S5.2); and the Elliptic result is
restated to include the fact that EOD moves the \emph{wrong} way (\S5.2).
\item \textbf{Text/table consistency.} We re-audited every numeric claim in the
paper against its cited table. Nine discrepancies were found and fixed --- the
five Reviewer~1 identified plus four more (Section~\ref{sec:extra}).
\item \textbf{DP scope.} The abstract now states the scope limitation in its own
sentence rather than in a parenthesis, and the positioning table (Table~1) now
carries a footnote so that our \checkmark\ in the DP column cannot be read as
equivalent to DP-FedAvg's record-level guarantee.
\end{itemize}

\subsection*{E5. Practical implications of requiring protected attributes at inference; trust-score limitations}

\begin{comment}
We also encourage you to expand the discussion of the practical implications of
requiring protected attributes at inference time and to make the limitations of
the composite trust score more explicit.
\end{comment}

Both done, and both expanded well beyond a sentence, since Reviewer~2 asked for
the same two things. See Sections~\ref{sec:r2-1} and~\ref{sec:r2-2} of this
letter: a new dedicated paragraph in \S6 (p.~32) analysing the lawful basis,
data-access, collection and purpose-limitation consequences for a deploying
institution, and a restructuring of \S6 so the trust score's caveats are stated
\emph{before} the score is presented.

\section{Response to Reviewer 1}

We thank Reviewer~1 for an exceptionally precise reading. The five text/table
mismatches in comment~6 were all real, and finding them required checking our
numbers against our own tables line by line; we are grateful for the effort.

\subsection*{R1.1 Softening the \texttt{robust\_bfwa} claim}
\label{sec:r1-1}

\begin{comment}
Section 5.8. The statement that robust-bfwa is the only aggregator that stays
strong under all three attacks is stronger than what Table 9 shows. The
trimmed-mean aggregator achieves comparable utility and even performs better
under the ALIE attack. I suggest softening this claim.
\end{comment}

The reviewer is correct and we have withdrawn the claim. Trimmed mean retains
$0.984/0.956/0.983$ AUC across the three attacks against our $0.992/0.953/0.990$,
so it survives the suite as we do and is indeed marginally better under ALIE. The
passage now reads \where{\S5.8, p.~26}:

\begin{newtextbox}
``$\dots$ FedAvg and plain BFWA collapse under the Gaussian attack ($0.51$ and
$0.48$ AUC) and Krum collapses under ALIE ($0.71$), whereas
\texttt{robust\_bfwa} ($0.992/0.953/0.990$) \emph{and the trimmed mean}
($0.984/0.956/0.983$) both stay strong under all three. We are careful here not
to overstate: the table does not support an exclusivity claim for our
aggregator, and the trimmed mean in fact retains marginally more utility than
ours under ALIE ($0.956$ vs.\ $0.953$). The accurate statement is that two of
the seven aggregators survive the full suite --- ours and the trimmed mean ---
and that among those two, ours is the one that also carries the fairness
constraint of Section~4.3. $\dots$ The honest summary is that
\texttt{robust\_bfwa} achieves a strong overall \emph{balance} $\dots$ rather
than dominating any metric.''
\end{newtextbox}

Checking this comment also revealed that the same overstatement had propagated
to the ablation discussion, where \texttt{robust\_bfwa} was again called ``the
only aggregator that stays strong across all three attacks''. That sentence is
corrected in the same way \where{\S5.3, p.~20}. While re-reading Table~9 we
further noticed that our list of cells where competitors beat us was incomplete:
trimmed mean attains the lowest EOD under fairness poisoning ($0.005$) and
multi-Krum the lowest EOD under Gaussian corruption ($0.017$). Both are now
named in the text.

\subsection*{R1.2 Sixteen versus fifteen baselines}

\begin{comment}
Number of baselines. The manuscript refers to sixteen baselines, whereas the
main comparison tables appear to contain fifteen. Please check whether one
baseline is missing from the tables or revise the reported number throughout the
manuscript.
\end{comment}

The reviewer is right: the tables contain fifteen baselines and the count of
sixteen was wrong. The discrepancy is \textbf{FedFB}, which we described among
the compared methods but never ran, because its FairBatch-style local group
reweighting coincides in our pipeline with the local soft-demographic-parity
objective already carried by FairFed and F$^2$GNN --- running it would have
duplicated an existing row rather than covering a new axis.

We have corrected the count to \textbf{fifteen} in all six places where it
appeared (abstract, \S1 contribution~5, \S1 relation-to-preliminary-version,
\S3.2, \S7 conclusion), removed FedFB from the list of \emph{evaluated}
baselines in \S3.2 and \S2, and retained it only in the related-work survey,
where it belongs as literature. To prevent the ambiguity recurring, \S3.2 now
ties the count directly to the tables \where{\S3.2, p.~15}:

\begin{newtextbox}
``The fifteen correspond exactly to the fifteen baseline rows of Tables~4--6,
above the rule that separates them from our own three variants; no method is
described here that does not appear there. (FedFB, surveyed in Section~2, is
\emph{not} among them: its FairBatch-style local group reweighting coincides
with the local soft-DPD objective already carried by FairFed and
F$^2$GNN in our pipeline, so running it would have duplicated an existing row
rather than added an axis.)''
\end{newtextbox}

\subsection*{R1.3 Uneven baseline coverage across datasets}

\begin{comment}
Baseline comparisons across datasets. The paper gives the impression that all
datasets are evaluated against the full set of baselines. However, the Credit
and especially the Elliptic datasets include fewer competing methods. This
should be stated explicitly in the experimental setup.
\end{comment}

Agreed --- this was left to be inferred from dashes in the tables, which is not
good enough. The experimental setup now carries a dedicated paragraph stating
the coverage exactly \where{\S3.2, p.~17}:

\begin{newtextbox}
``\textbf{Baseline coverage is not uniform across datasets.} $\dots$ The
\emph{full} fifteen-baseline comparison is run only on the three datasets that
combine a genuine protected attribute with a tractable cost: \textbf{German},
\textbf{Bail}, and \textbf{Pokec}. \textbf{Credit} ($30{,}000$ nodes) is run
against the seven core baselines $\dots$ plus DP-FedAvg, but not against the
five 2025--2026 reimplementations or FaVGNN/FDP-Fair. \textbf{Elliptic}
($203{,}769$ nodes) is run against only FedAvg-GAT and FairSIN, the two
baselines that scale to it within our CPU budget $\dots$ Consequently the Credit
and Elliptic columns of Tables~4--6 should be read as comparisons against
\emph{that} restricted set, and a `--' always means \emph{not run}, never a
failed or omitted result.''
\end{newtextbox}

The consequence is also propagated to the places where those two datasets are
discussed: the Elliptic paragraph in \S5.2 now says explicitly that only two
baselines were run at that scale, so the result is ``not a broad competitive
ranking''; the introduction flags the restriction when it announces the
evaluation; and a new limitation item (4b) records it.

\subsection*{R1.4 Single-seed robustness experiments}

\begin{comment}
Robustness experiments. Since robustness is one of the main contributions, the
conclusions are currently based on single-seed experiments. Adding a few
additional runs would strengthen the evidence. If this is not feasible, the
claims should be slightly moderated.
\end{comment}

We have taken the second option and moderated the claims, and we want to be
transparent about why. A multi-seed replication of the $7\times3$
aggregator--attack grid is a substantial compute commitment that we could not
complete within the revision window without compromising the rest of the
response; we would rather state the limit of the present evidence precisely than
add two seeds and still be unable to support a ranking.

What we have done instead is to bound the claims to what one seed can carry. A
new paragraph \where{\S5.8, p.~26} separates the order-of-magnitude findings we
rely on from the small differences we explicitly decline to interpret:

\begin{newtextbox}
``\textbf{How far this single-seed evidence reaches.} These are
\emph{single-seed} runs on Bail, and we deliberately do not lean on them harder
than that allows. Differences of a few thousandths in AUC or EOD between the
surviving aggregators --- ours, the trimmed mean, multi-Krum --- are well within
the range seed noise alone produces elsewhere in this paper (TrustFedGNN's own
Bail DPD has a standard deviation of $0.013$ over five seeds), so we make no
ranking claim among them, and no conclusion in this paper rests on such a gap.
What a single seed \emph{does} support, because these effects are
order-of-magnitude rather than marginal, are the two qualitative findings we
actually rely on: that several aggregators \emph{collapse} outright under a
given attack (FedAvg and plain BFWA to $\approx 0.5$ AUC under Gaussian noise,
Krum to $0.71$ under ALIE), and that plain BFWA's fairness is captured by the
fairness-poisoning attack --- an EOD of $0.126$, more than five times the
$0.024$ of the screened variant. Replicating the full $7\times3$
aggregator--attack grid over multiple seeds, with error bars and with adaptive
attackers who know the screening rule, is the natural next step $\dots$''
\end{newtextbox}

The limitations section was rewritten to match, and now names the robustness
study as the place where statistical power runs out and identifies its
multi-seed replication as ``the most valuable single addition to this
evaluation'' \where{\S7, Limitation (4), p.~34}.

\subsection*{R1.5 The energy column}

\begin{comment}
Energy consumption results. The energy column in Table 11 reports 0.00 Wh for
all methods, making it difficult to draw any comparison. Either report the
values with higher precision or clarify that the measurements fall below the
resolution of the measurement tool.
\end{comment}

Investigating this comment showed the problem was worse than a precision issue,
and we thank the reviewer for prompting the check. Our energy proxy is
wall-clock $\times$ an assumed CPU draw; wall-clock was \emph{never logged} for
those canonical small-graph runs, so the proxy silently received zero and the
column printed ``0.00'' for every method. It was not a measurement below
resolution --- it was not a measurement at all. Reporting it as though it were a
comparison would have been misleading.

We therefore made two changes rather than adjusting the precision:

\begin{enumerate}[leftmargin=1.4em,itemsep=2pt]
\item \textbf{The energy column is removed from Table~11}, and the caption now
says why, so no reader infers energy parity from a row of zeros:
\begin{newtextbox}
``Energy is \emph{not} reported here: wall-clock was not logged for these
canonical small-graph runs, so the proxy of Section~4.6 has no input for them.
We report it instead where the measurement exists, on the large-scale study
(Table~10), rather than print an uninformative zero.''
\end{newtextbox}
\item \textbf{Energy is now reported where the measurement does exist.} The
large-scale ogbn-products study logged per-run wall-clock, so Table~10 gains an
\textbf{Energy (Wh)} column derived from it, and the numbers are informative:
FedAvg-GAT $26.2$\,Wh, FairSIN $24.3$\,Wh, DP-FedAvg $22.6$\,Wh, FaVGNN
$50.4$\,Wh, TrustFedGNN $66.2$\,Wh. \S5.10 now reads the comparison out
explicitly --- roughly $40$\,Wh, about $16$\,gCO$_2$e, as the one-off cost of
training a fair, private and robust model on a $2.4$M-node graph --- while
labelling the figure a proxy on fixed hardware rather than a metered
measurement.
\end{enumerate}

No experiment was re-run for this: the wall-clock values were already in
Table~10 and the energy column is an arithmetic transformation of them. We also
patched the two table-generating scripts so that an unmeasured run now emits
``--'' instead of a spurious zero, preventing the defect from recurring.

\subsection*{R1.6 Text/table inconsistencies}
\label{sec:r1-6}

\begin{comment}
I also noticed a few inconsistencies between the text and tables that should be
corrected $\dots$
\end{comment}

All five were real. Each is corrected as follows.

\textbf{(a) FedFACT DPD on Bail.} The text said $0.015$; Table~5 says $0.021$.
The table is correct and the text now reads $0.021$ \where{\S5.2, p.~19}. (The
surrounding claim --- that FedFACT achieves the lowest DPD among non-private
methods on Bail --- remains true at the corrected value.)

\textbf{(b) FairFed $p$-value on German.} The text quoted $p{=}0.006$ while Table~3 shows $0.01$. The table rounds $p$-values to two decimals; the text now
quotes the rounded value $p{=}0.01$, and the caption of Table~3 states the
rounding convention explicitly so that text and table cannot drift again.

\textbf{(c) The AUC claim does not hold for FairGNN.} Correct, and this was an
overgeneralisation on our part. FairGNN's $p{=}0.16$ is not significant. The
sentence now names both the rule and the exception \where{\S5.2, p.~19}:

\begin{newtextbox}
``$\dots$ on German its AUC is significantly below \emph{most} non-private
baselines --- FedAvg-GAT, FairSIN, F$^2$GNN ($p{=}0.02$) and FairFed
($p{=}0.01$) --- the cost of DP plus the fairness/robustness constraints. The one
non-private baseline it is \emph{not} significantly below is FairGNN
($p{=}0.16$), whose adversarial objective is itself unstable on a graph this
small (AUC $0.668\pm0.084$); we state this exception rather than generalise over
it.''
\end{newtextbox}

\textbf{(d) The attribute-inference baseline: $0.50$ versus $0.53$.} This was a
genuine defect in the figure, not merely inconsistent phrasing, and we are
grateful it was caught. Our attack samples targets \emph{in equal numbers from
each true group} ($2{,}000$ per group), so an adversary that ignores the release
entirely cannot beat $0.500$ on that target set --- which is the baseline the
table reports. The figure, however, drew a line labelled ``chance (base rate):
0.53'', which is the majority-group \emph{prevalence} in Bail, the correct chance
level only for an \emph{unbalanced} target set. The label was wrong.

We have re-rendered Figure~5 from the already-logged attack results (no
re-training; the seven $(\epsilon,\text{accuracy})$ pairs are exactly those in
Table~8, so figure and table are now guaranteed to agree). The corrected figure
draws the chance line at $0.500$ and labels the $0.53$ prevalence line for what
it is. The text \where{\S5.5, p.~24} and both captions now explain the
distinction, and the plotting script was patched so the mislabelling cannot
return.

We want to be explicit that we resolved this in the direction that makes our own
claim \emph{harder} to support, not easier. Reporting against $0.53$ would have
been superficially flattering --- every measured accuracy ($0.502$--$0.513$) lies
\emph{below} it --- but it would have compared a balanced-target accuracy against
an unbalanced-target baseline, and an attack that appears to perform below chance
signals a misspecified evaluation rather than a privacy guarantee. The balanced
$0.500$ requires the mechanism to beat a coin flip rather than a class-prior
guess, which is the test we should be held to.

Re-examining the comparison this way also let us make a stronger and more
defensible statement than before. Each row of Table~8 is a Monte-Carlo estimate
over $4{,}000$ target draws, so its standard error is
$\sqrt{0.25/4000}\approx0.008$; the largest excess over chance in the entire
sweep is $0.013$, or $1.6$ standard errors. We therefore added:

\begin{newtextbox}
``Measured against that harder reference, the residual leakage is not merely
small but statistically undetectable at our sample size $\dots$ \emph{At no
privacy budget is the attack's accuracy significantly different from $0.500$}
(two-sided $p\ge0.10$ throughout). We therefore claim that FTGD reduces the
differencing adversary to chance, not merely close to it --- while noting that
this bounds leakage only at the resolution $4{,}000$ trials can resolve, so an
advantage below roughly one percentage point would not be visible to this test.''
\end{newtextbox}

\textbf{(e) German DPD in the ablation versus the main table.} The ablation
reports $0.045$ and Table~5 reports $0.040\pm0.035$. The reason is that the
ablation is a \emph{separate set of runs} with its own matched seed set --- so
that all four configurations are compared on identical seeds --- whereas the main
tables aggregate the full seed sets (10 on German, 5 on Bail). We had not stated
this, which made the two look like they should agree exactly. The ablation
caption now explains it and shows that all four affected cells differ by less
than one standard deviation of the multi-seed mean \where{Table~7, p.~23}.

\subsection*{R1.7 Scope of the DP guarantee in the abstract}
\label{sec:r1-7}

\begin{comment}
Finally, the abstract and the description of the differential privacy guarantee
should make it clearer that the guarantee applies to the released fairness
statistics, rather than the entire model update.
\end{comment}

Agreed. The scope was previously signalled in the abstract only by a
parenthetical (``we state precisely what this guarantee covers, and what it does
not, in the text''), which is too easy to skim past for a limitation this
important. It is now a sentence of its own:

\begin{newtextbox}
``We emphasise at the outset that this $(\epsilon,\delta)$ guarantee covers the
\emph{released fairness statistic only}, and \emph{not} the full transmitted
model update, which retains sensitive-attribute dependence through FSER and the
fairness-gradient masks; the comparison against full-gradient DP-SGD is
therefore not like-for-like, and we make the exact scope precise in
Section~4.2.''
\end{newtextbox}

Beyond the abstract, we made three further scope-tightening changes. The
abstract's closing claim and the conclusion now say ``\emph{statistic-level}
sensitive-attribute DP'' rather than ``sensitive-attribute DP''. Contribution~2
in \S1 already carried the caveat and is unchanged. Most importantly,
\textbf{Table~1} previously gave TrustFedGNN a plain \checkmark\ in the DP
column, directly beside DP-FedAvg, PUFFLE and FDP-Fair, which implied
guarantees of equal scope; that entry now carries a marker and the caption
explains:

\begin{newtextbox}
``$^{\ddagger}$Our DP entry is \emph{not} equivalent to the record-level
guarantees of DP-FedAvg, PUFFLE, FDP-Fair or FaVGNN: TrustFedGNN privatises the
\emph{released fairness statistic} rather than the whole model update
(Section~4.2), so this column marks the presence of a formal
$(\epsilon,\delta)$ guarantee, not guarantees of equal scope.''
\end{newtextbox}

\section{Response to Reviewer 2}

We thank Reviewer~2 for identifying the two places where the paper's framing,
rather than its content, was doing readers a disservice. Both comments led to
substantive rewrites.

\subsection*{R2.1 Practical implications of requiring the protected attribute at inference}
\label{sec:r2-1}

\begin{comment}
FSER requires the protected attribute at inference time. This is currently
called a ``governance nuance'' in one line. Given the paper's emphasis on EU AI
Act compliance, this deserves a dedicated paragraph in Section 6 if possible, by
the authors to explain what this means practically for a deploying institution,
such as what data access agreements or any legal obligations it creates.
\end{comment}

We agree entirely: ``governance nuance'' understated what is the single design
decision in TrustFedGNN with the largest practical and legal footprint. \S6 now
carries a dedicated paragraph \where{\S6, p.~32} that first makes the technical
requirement concrete --- because the model attends over a neighbourhood, scoring
one applicant requires the protected attribute of the \emph{counterparties} in
that applicant's neighbourhood, not only of the applicant --- and then works
through four consequences for a deployer:

\begin{newtextbox}
\textbf{First, lawful basis.} Gender, race and ethnicity are special-category
data under GDPR Art.~9; processing them in an automated decision requires an
Art.~9(2) condition, and in several member states the substantial-public-interest
route additionally needs a national-law basis. The EU~AI~Act anticipates this
tension: Art.~10(5) permits processing special-category data \emph{for the
purpose of bias detection and correction} in high-risk systems, subject to
safeguards including technical limits on re-use, pseudonymisation, and deletion
once bias correction is complete. FSER is squarely a bias-correction mechanism
and so is the kind of processing that provision contemplates --- but an
institution must document that reliance rather than assume it, and Art.~10(5) is
a permission to process, not a waiver of Art.~9 or of a DPIA.\\[0.35em]
\textbf{Second, data-access agreements.} In the cross-silo setting the
counterparties in a neighbourhood may be another institution's customers, so
inference-time access to $s$ can require inter-institutional agreements that
plain FedAvg does not --- a governance cost that partly offsets the ``no raw data
leaves the silo'' benefit of federation.\\[0.35em]
\textbf{Third, collection risk.} Where $s$ is not already held, obtaining it
means either new collection (consent, minimisation, retention limits) or
\emph{inference} of a protected attribute from proxies, which is itself a
discriminatory-profiling risk and which we do not recommend.\\[0.35em]
\textbf{Fourth, purpose limitation.} $s$ must reach the attention mechanism
without reaching the classifier as a feature: our implementation keeps it out of
$\mathbf{X}$, but that separation is an engineering invariant a deployer must
verify and an auditor must test, not a property the architecture enforces by
itself.
\end{newtextbox}

The paragraph closes by telling a reader who \emph{cannot} meet these
obligations what to do: FSER is then the wrong component, and the ablation
(Table~7) shows FTGD and BFWA operating without it at a quantified fairness cost.
We also added a new limitation item (12) recording that ``until this is done,
TrustFedGNN is only deployable where such processing is lawful and governed'',
and a future-work direction on removing the dependence altogether --- by
distilling FSER into an $s$-free student model, or by driving the edge correction
from the DP-noised statistics instead of a hard group indicator.

\subsection*{R2.2 Making the trust-score limitation more prominent}
\label{sec:r2-2}

\begin{comment}
Also, the paper mentions the composite trust score is not decisive evidence, but
this appears after the score is already presented. If possible, highlight this
limitation more prominently, which can be helpful for readers.
\end{comment}

A fair criticism: a caveat that arrives after the number has been read has
already failed. We have inverted the order. The subsection is retitled
``\textbf{Composite trust score: read this caveat first}'' and now opens with the
limitation, before the score is defined \where{\S6, p.~29}:

\begin{newtextbox}
``Because a single headline number invites exactly the over-reading we want to
avoid, we state its limits \emph{before} presenting it rather than after. The
composite score below is a \emph{communication tool} $\dots$ It is not a
certification of trustworthiness, not a substitute for the per-axis evidence,
and not decisive evidence of superiority. Its ranking depends on choices we make
--- the reference constants, the unit weights, and the power $p$ --- and a reader
who chooses different constants may obtain a different ordering. Auditors,
deployers and reviewers should therefore treat the per-axis results (Tables~4--11)
as the primary empirical findings, and the composite as a summary of them. With
that framing fixed, the construction is as follows.''
\end{newtextbox}

The corresponding limitation item was also strengthened, and now adds a point we
had not made anywhere: the composite covers only the axes computable from a given
run, so the three-axis version in Table~12 omits calibration and robustness
entirely, and ``a high score is a statement about the axes included, not about
the system as a whole'' \where{\S7, Limitation (10), p.~35}.

\section{Further corrections made on our own initiative}
\label{sec:extra}

Prompted by Reviewer~1's comment~6, we re-audited every numerical claim in the
paper against the table it cites. Four further discrepancies surfaced, none
raised by the reviewers; we list them for completeness.

\begin{enumerate}[leftmargin=1.5em,itemsep=3pt]
\item \textbf{A limitation contradicted the paper.} Limitation (4) still read
``we $\dots$ do not run significance tests'', which was true of an earlier draft
but false once the paired Wilcoxon tests of Table~3 were added. The item is
rewritten to say precisely where the paper \emph{is} statistically supported
(German, Bail, Pokec at $\geq 5$ seeds) and where it is not (Credit/Elliptic at
2--3 seeds; robustness, large-scale and calibration at one seed).

\item \textbf{An overstated claim about PUFFLE.} \S5.2 said PUFFLE ``reaches the
lowest DPD/EOD of \emph{any} method on Bail ($0.001$)''. Its DPD of $0.001$ is
\emph{tied} with FDP-Fair and DP-FedAvg, and its EOD of $0.006$ is beaten by
TrustFedGNN-Robust ($0.005$). Corrected to a tie among the full-gradient-DP
methods.

\item \textbf{An imprecise margin.} The AUC margin over the privacy-preserving
methods on Bail was quoted as ``$+0.40$''; the three actual margins are $+0.38$,
$+0.40$ and $+0.44$. Now reported as a range.

\item \textbf{A wall-clock ratio outside its stated range.} The TrustFedGNN /
FedAvg per-round cost was quoted as ``$2.4$--$2.5\times$''; Table~10 gives
$2.45\times$ and $2.53\times$. Corrected to $2.45$--$2.53\times$.
\end{enumerate}

Additionally, and in the same spirit as R1.1, we restated the Elliptic result in
\S5.2. It previously read that TrustFedGNN ``attains strong subgroup fairness at
detection accuracy comparable to the non-fair backbone'', which quietly omitted
that EOD moves the \emph{wrong} way on that dataset ($0.020 \to 0.075$) --- a
tension our own limitations section already acknowledged. The passage now gives
DPD, AUC \emph{and} the adverse EOD result together.

\section{Closing}

We believe the manuscript is substantially more accurate for these comments. The
technical content is unchanged --- no result has been altered, and no new
experiment was needed --- but several claims are now bounded by what our
evidence actually supports, the text and tables agree everywhere we could check
them, the scope of the privacy guarantee is unambiguous, and the two governance
questions the Editor and Reviewer~2 raised (inference-time access to protected
attributes, and the standing of the composite trust score) are treated at the
length they deserve.

We would like to thank the Reviewers once more for the care they invested, and
the Editor for the opportunity to revise.

\vspace{1em}
Sincerely,\\[0.4em]
Minh Ngoc Dinh, on behalf of all authors\\
\textit{Corresponding author}

\end{document}
